\chapter{System Analysis}

\section{Functional Requirements}
To achieve the project objectives, Vigil-Core must fulfill the following functional requirements:
\begin{itemize}
    \item \textbf{File Upload and Validation:} The system must allow users to securely upload potential malware samples and network capture (PCAP) files for analysis.
    \item \textbf{Automated Pipeline:} The backend must automatically route uploaded files through a 5-phase pipeline: Dump Validator $\rightarrow$ Region Scanner $\rightarrow$ Payload Extractor $\rightarrow$ IOC Extractor $\rightarrow$ Report Generator.
    \item \textbf{YARA and Hashing Integration:} The platform must run uploaded samples against predefined YARA rules and perform fuzzy hashing (using pydeep/ssdeep) to detect anomalies and variants.
    \item \textbf{Data Visualization:} The UI must dynamically render charts (e.g., using Recharts) to visualize entropy, protocol distributions, and YARA match statistics.
    \item \textbf{User Management:} The system must support user authentication, allowing analysts to maintain a history of their scans and reports.
\end{itemize}

\section{Non-Functional Requirements}
\begin{itemize}
    \item \textbf{Performance and Scalability:} The system must handle large files gracefully. Specifically, it must process memory dumps up to 8GB without crashing, utilizing chunked processing techniques. The UI must maintain a low latency (target $\sim$4.3ms) to ensure a smooth user experience.
    \item \textbf{Cross-Platform Compatibility:} While initially developed on Windows, the backend must be fully functional on Linux environments, requiring the correct installation and configuration of dependencies like \texttt{libpcre3}, \texttt{libpcap}, and \texttt{python3-dev}.
    \item \textbf{Reliability:} API responses from the backend to the frontend must be structured and stable, with comprehensive error handling to manage malformed files or unsupported formats.
\end{itemize}

\section{Feasibility Study}
\subsection{Technical Feasibility}
The project relies on well-documented and widely supported technologies: Python for the backend scripting, Flask/FastAPI for API routing, and React for the frontend. The integration of C-based libraries (like YARA and libpcap) via Python bindings is technically sound but requires meticulous environment configuration, which was successfully achieved during the testing phase.

\subsection{Operational Feasibility}
Vigil-Core is designed to be deployed in standard SOC environments. By containerizing the application or providing clear dependency installation scripts for Linux, it operates seamlessly on standard analyst workstations. The intuitive UI ensures that even junior analysts can quickly utilize the platform's capabilities without extensive training on command-line utilities.
